
\section{Related Work}
\label{sec:relatedwork}
%\subsection{Data Augmentation}
% \noindent\textbf{Graph Augmentation.} Data augmentation is an effective technique to alleviate the overfitting problem in training deep networks~\citepp{wang2019implicit}. While in the framework of contrastive learning, Data augmentation is to produce two views which are semantically similar~\citep{he2020momentum, chen2020simple}. The standard graph augmentation are attribute masking, edge permutation and node dropping. Being formally studied in~\citep{you2020graph}, these graph augmentation are considered to impose certain human prior to the SSL model (\ie,Edge perturbation means semantic robustness against connectivity variations). In addition to this, Some prior works have also investigated adaptive graph augmentations~\citep{zhu2021graph}, using heuristics such as node centrality or PageRank centrality~\citep{page1999pagerank} to mask different edges with different probabilities. This improves the quality of the augmented graphs by helping these transformations preserve semantic similarity (\eg, by making it less likely to mask critical edges that connect otherwise disjoint parts of the graph). Noted that graph augmentation conduct directly in the input space, which will, after graph encoder, introduce large bias to the node features in hidden features space.

% \subsubsection{Input Space Augmentation}
%\vspace{0.05cm}
\noindent\textbf{Data Augmentation.} %Broadly studied in various domains, 
Augmentations are usually performed in the input space. In computer vision, image transformations, \ie, rotation, flipping, color jitters, translation,  noise injection~\citep{shorten2019survey}, cutout and random erasure~\citep{devries2017dataset} are  popular. In neural language processing, %input space augmentations include
token-level random augmentations, \eg, synonym replacement, word swapping, word insertion, and deletion~\citep{wei2019eda} are used. In transportation, conditional augmentation of road junctions is used \cite{coltrane}. 
%
In the graph domain, %input space augmentation equals graph augmentation: 
attribute masking, edge permutation, and node dropout are popular \citep{you2020graph}. Sun \etal \cite{uai_ke} use adversarial graph perturbations. Zhu \etal \citep{zhu2021graph} use adaptive graph augmentations based on %heuristics such as 
the node/PageRank centrality~\citep{page1999pagerank} to mask  edges with varying probability.
%It should be emphasized that, unlike other input space DA, graph augmentation will introduce bias to the hidden feature space.
%
%Although it has been empirically proven that data augmentation is  effective for self-supervised graph learning, it is not as efficient for learnable augmentation approaches on huge graphs. 

%\subsection{Feature Augmentation}
\vspace{0.1cm}
\noindent\textbf{Feature Augmentation.} Samples can be augmented in the feature space instead of the input space~\citep{feng2021survey}. Wang \etal \citep{wang2019implicit} augment the hidden space features, resulting in auxiliary samples with the same class identity but different semantics. A so-called channel augmentation  perturbs the channels of feature maps~\citep{wang2019implicit} while  GCL approach, COSTA \cite{zhang2022costa}, augments features via random projections. Some few-shot learning approaches augment features  \cite{Zhang_2022_CVPR} while others  estimate the ``analogy'' transformations between samples of known classes to apply them on samples of novel classes \citep{hariharan2017low,schwartz2018delta} or mix foregrounds and backgrounds \cite{zhang2019few}. 
% \citep{paschali2019data} use iterative affine transformations and projections to maximally "stretch" an example along the class-manifold.
%While feature space augmentations are popular in semi-supervised learning, one-shot, and few-shot learning, 
However, ``analogy'' augmentations are not applicable to contrastive learning due to the lack of labels.  %Recently, Zhang \etal ~\cite{zhang2022costa} proposed COSTA, a state-of-the-art GCL feature augmentation  via random projections. %It achieve SOTA result on graph contrastive learning in several datasets.
%some work combine feature augmentation with contrastive learning, 

%Representation
%\subsection{Graph Contrastive Learning}
\vspace{0.1cm}
\noindent\textbf{Graph Contrastive Learning.} 
CL is popular in computer vision, NLP~\citep{he2020momentum, chen2020simple, gao2021simcse}, and graph learning. %Compared with the vision domain, GCL uses different augmentation and contrasting strategies. 
In the vision domain,  views are formed by augmentations at the pixel level, whereas in the graph domain, data augmentation may act on node attributes or the graph edges. %Similarly, in the contrasting mode,
%For the contrastive setting, 
GCL often explores node-node, node-graph, and graph-graph relations for contrastive loss which is similar to contrastive losses in computer vision. Inspired by SimCLR~\citep{chen2020simple}, GRACE~\citep{zhu2020deep}  correlates graph views by pushing closer representations of the same node in different views and separating representations of different nodes, and  Barlow Twin~\citep{zbontar2021barlow}   avoids the so-called dimensional collapse %(embedding vectors occupy a lower-dimensional subspace than their dimension)
~\citep{jing2021understanding}. % but are incompatible with other contrastive designs. %Thus, we propose a loss-independent feature augmentation which does not suffer from the dimension collapse, and improves the performance of InfoNCE and Barlow twin. 

In contrast, we study spectral feature augmentations to perturb/rebalance singular values of both views. We outperform  feature augmentations such as COSTA~\cite{zhang2022costa}.

% By adapting DeepInfoMax~\citep{bachman2019learning} to graph representation learning, %DGI~\citep{velickovic2019deep} with GCNs encodes the input graph. %, and the graph embeddings are obtained via applying a readout function. % to the node embeddings. 
% DGI~\citep{velickovic2019deep} learns the embedding of nodes and graphs via maximizing their mutual information, which discriminates between nodes of the original graph and corrupted graphs. Inspired by the SimCLR~\citep{chen2020simple}, GRACE~\citep{zhu2020deep}  correlates graph views by pushing closer representations of the same node in different views and pushing apart representations of different nodes. Another example of  SimCLR strategy is the recent GraphCL method~\citep{hafidi2020graphcl}. In contrast to GRACE, which learns node embedding, GraphCL learns embeddings for graph-level tasks.